\documentclass[11pt]{article}
\usepackage{graphicx}
\usepackage[T1]{fontenc}
\usepackage[polish]{babel}
\usepackage{tabularx}
\usepackage[table,xcdraw]{xcolor}
\usepackage[utf8]{inputenc}
\usepackage{lmodern}
\usepackage{multirow}
\usepackage{array}
\usepackage{booktabs}
\selectlanguage{polish}
\usepackage{titlesec}
\usepackage{amsmath}
\usepackage{esint}
\usepackage{textpos}
\usepackage{chngpage}
\usepackage{calc}
\usepackage{placeins}
\usepackage{longtable}
\usepackage{hyperref}

\setlength\parindent{24pt}

\begin{document}

	
\begin{titlepage}				
\centering

{\Large{Wydział Matematyki i Nauk Informacyjnych Politechniki Warszawskiej}}

\vspace{1cm}		
{\includegraphics[scale=0.25]{logo}}
\vspace{3cm}

\Huge\textbf{Dokumentacja} 
\bigbreak

{\Huge Analiza ilości terenów zielonych na terenie dzielnic Warszawy}
\vspace{1cm}

{\Large Bartosz Jasiński, Marta Rutkowska}
\vspace{1cm}

\Large{v1.0}
\vspace{1cm}

{\itshape {\Large \today}}

\vfill		

\end{titlepage}


\tableofcontents


\newpage

	\section{Wstęp}
	\subsection{Realizowany cel biznesowy}
		Analiza stosunku pola przestrzeni zielonych do ilości mieszkańców, procentu zazielenia i gęstości zaludnienia dzielnic Warszawy na podstawie zdjęć satelitarnych oraz danych administracyjnych i demograficznych Warszawy.
		
	\subsection{Użytkowanie aplikacji}
		Aplikacja wyświetla mapę Warszawy podzieloną na dzielnice na które będzie nałożony pół-przejrzysty kolor (od czerwonego do zielonego) w zależności od procentu powierzchni zajmowanej przez tereny zielone w danej dzielnicy. Użytkownik ma możliwosć wybrania z mapy interesującej go dzielnicy lub zbliżenia środkowym przyciskiem myszki interesującego go obszaru i odświeżenia informacji na temat zbliżonego obszaru. Podane zostaną następujące informacje:
    \begin{itemize}
      \item gęstości zaludnienia
      \item liczby mieszkańców  
      \item ilości powierzchni terenów zielonych przypadających w danym obszarze na mieszkańca
      \item jaki procent powierzchni danego obszaru zajmują obszary zielone
    \end{itemize}
Zapytania w zależności od rodzaju będą przeprowadzane w czasie rzeczywistym (podczas zbliżania mapy) lub będą zwracane wcześniej wyliczone wartości (w czasie klikania na daną dzielnicę). 


	\subsection{Technologia}
		Projekt wykonany w technologii: \textbf{Python} po stronie serwera najprawdopodobniej z użyciem frameworku \textbf{Tornado}, razem z bibliotekami pomocnymi w obsłudze plików \textbf{csv}, \textbf{JavaScrip} oraz \textbf{HTML} po stronie przeglądarki.  

	
	\section{Dane aplikacji}
	\subsection{Źródła danych}
		Dane użyte w projekcie pochodzą z publicznie udostępnianych, darmowych źródeł wymienionych poniżej.
	
	\begin{itemize}
	\item Zdjęcia satelitarne z Google Maps API: https://maps.googleapis.com
	\item Dane na temat podziału miasta Warszawa na dzielnice pozyskane ze strony http://www.gugik.gov.pl, pełny link do danych: http://www.gugik.gov.pl/geodezja-i-kartografia/pzgik/dane-bez-oplat/dane-z-panstwowego-rejestru-granic-i-powierzchni-jednostek-podzialow-terytorialnych-kraju-prg
	\item Dane na temat ludności dzielnic Warszawy, stan na 2017 rok z portalu http://stat.gov.pl, pełny link do danych: http://stat.gov.pl/obszary-tematyczne/ludnosc/ludnosc/powierzchnia-i-ludnosc-w-przekroju-terytorialnym-w-2017-r-,7,14.html
	\end{itemize}

	\subsection{Baza danych aplikacji}
	Dane na temat ludnosci poszczególnych dzielnic Warszawy przechowywane są w pliku csv po stronie serwera. Podział miasta na dzielnice wczytywany jest z pliku Shapefile.

	\section{Utrzymanie aplikacji}
		Utrzymanie aplikacji będzie opierało się głównie o odswieżanie danych na temat ludnosći dzielnic Warszawy oraz ewentualnym odswieżeniu danych o granicach administracyjnych dzielnic w razie ich zmiany, lub utworzenia nowych jednostek oraz uaktualniania zdjęć satelitarnych które służą do szacowania ilości terenów zielonych.


\end{document}
